% ---------------------------------------------------------------------------
% Apna Hisab - User Guide (Admin & Shop)
% Compile with:  pdflatex USER_GUIDE.tex   (run twice for the table of contents)
% Needs a standard TeX Live / MiKTeX install - no exotic packages.
% ---------------------------------------------------------------------------
\documentclass[11pt,a4paper]{article}

\usepackage[T1]{fontenc}
\usepackage[utf8]{inputenc}
\usepackage[a4paper,margin=1in]{geometry}
\usepackage{xcolor}
\usepackage{enumitem}
\usepackage{tabularx}
\usepackage{booktabs}
\usepackage{titlesec}
\usepackage{fancyhdr}
\usepackage{parskip}
\usepackage[hidelinks]{hyperref}

% ---- Colours ----
\definecolor{brand}{HTML}{4F39C6}
\definecolor{ink}{HTML}{1F2937}
\definecolor{muted}{HTML}{6B7280}
\definecolor{soft}{HTML}{F3F4F6}

\hypersetup{
  colorlinks=true,
  linkcolor=brand,
  urlcolor=brand,
  pdftitle={Apna Hisab - User Guide},
  pdfauthor={Sarathi Technologies}
}

% ---- Section styling ----
\titleformat{\section}{\Large\bfseries\color{brand}}{\thesection}{0.6em}{}
\titleformat{\subsection}{\large\bfseries\color{ink}}{\thesubsection}{0.6em}{}
\titleformat{\subsubsection}{\normalsize\bfseries\color{ink}}{\thesubsubsection}{0.6em}{}

% ---- Header / footer ----
\pagestyle{fancy}
\fancyhf{}
\renewcommand{\headrulewidth}{0.4pt}
\lhead{\small\color{muted}Apna Hisab}
\rhead{\small\color{muted}User Guide}
\cfoot{\small\color{muted}\thepage}

% ---- Helpers ----
\newcommand{\role}[1]{\textbf{\color{brand}#1}}
\newcommand{\menu}[1]{\textsf{#1}}
\newcommand{\key}[1]{\fbox{\small\texttt{#1}}}

\newenvironment{note}
  {\par\medskip\noindent\begin{tabular}{@{}p{0.03\textwidth}p{0.9\textwidth}@{}}
   \textbf{!} & \itshape}
  {\end{tabular}\par\medskip}

% ---- Table column type ----
\newcolumntype{L}{>{\raggedright\arraybackslash}X}

\setlist[itemize]{leftmargin=1.4em,itemsep=2pt,topsep=3pt}
\setlist[enumerate]{leftmargin=1.6em,itemsep=3pt,topsep=3pt}

% ===========================================================================
\begin{document}

% ---- Title page ----
\begin{titlepage}
\centering
\vspace*{3cm}
{\Huge\bfseries\color{brand} Apna Hisab\par}
\vspace{0.4cm}
{\Large Billing \& Accounting for Your Shop\par}
\vspace{1.2cm}
{\LARGE\bfseries User Guide\par}
\vspace{0.3cm}
{\large For Administrators and Shop Owners\par}
\vspace{2cm}
\begin{tabular}{@{}ll@{}}
\textbf{Application} & Apna Hisab \\
\textbf{Shop app}    & \url{https://apnahisab.sarathitechnologies.com/shop/login} \\
\textbf{Admin panel} & \url{https://apnahisab.sarathitechnologies.com/console/admin/login} \\
\end{tabular}
\vfill
{\color{muted}\small Published by Sarathi Technologies \textbullet{} Revision 1.0\par}
\end{titlepage}

\tableofcontents
\clearpage

% ===========================================================================
\section{Introduction}

Apna Hisab is a web-based billing and accounting system for Indian shops. It
handles sales and purchase invoices, GST, inventory, customer and supplier
credit, expenses, and financial reports --- all from a web browser, with nothing
to install.

The system has \textbf{two separate panels} that share one web address:

\begin{center}
\begin{tabularx}{\textwidth}{@{}p{0.26\textwidth}L@{}}
\toprule
\textbf{Panel} & \textbf{Who uses it and what for} \\
\midrule
\role{Admin panel} (operator console) &
The platform administrator onboards shops, manages subscriptions and plans,
and oversees users. Opens under \texttt{/console}. \\[4pt]
\role{Shop app} &
Shop owners and their staff do daily billing, inventory, payments and reports.
Opens under \texttt{/shop}. \\
\bottomrule
\end{tabularx}
\end{center}

\begin{note}
In a single browser you can be signed into \emph{either} the admin panel or a
shop, not both at the same time (they share one web address). If you need both
open together, use a second browser or a private/incognito window.
\end{note}

\subsection{Addresses at a glance}

\begin{center}
\begin{tabularx}{\textwidth}{@{}p{0.22\textwidth}L@{}}
\toprule
\textbf{You are\ldots} & \textbf{Open this address} \\
\midrule
A shop owner / staff & \url{https://apnahisab.sarathitechnologies.com/shop/login} \\[3pt]
The administrator    & \url{https://apnahisab.sarathitechnologies.com/console/admin/login} \\
\bottomrule
\end{tabularx}
\end{center}

\subsection{Signing in and out}
\begin{itemize}
  \item Enter your email and password, then select \menu{Sign in}.
  \item Your session is kept securely; you stay signed in until you sign out or
        it expires.
  \item To sign out, open your profile menu (top-right) and choose \menu{Sign out}.
  \item If you forget a password, ask your administrator (for shop staff) or the
        platform administrator (for shop owners) to reset it.
\end{itemize}

% ===========================================================================
\clearpage
\section{Part A --- Administrator Guide}

The administrator runs the platform. This section is for the person who holds the
\role{Platform Admin} account.

\subsection{Opening the admin panel}
Go to \url{https://apnahisab.sarathitechnologies.com/console/admin/login} and sign
in with your administrator credentials. You land on the platform dashboard.

\begin{note}
Change the default administrator password immediately after your first sign-in,
from \menu{Admin \textrightarrow{} Settings} (or your profile menu).
\end{note}

\subsection{Onboarding a new shop}
This is how a shop and its owner login are created.
\begin{enumerate}
  \item Open \menu{Platform \textrightarrow{} Businesses}.
  \item Select \menu{Onboard business} (or \menu{Add business}).
  \item Fill in the shop details: business name, owner name, owner email,
        phone, and a starting password for the owner.
  \item Choose a subscription plan for the shop (see \S\ref{sec:plans}).
  \item Save. The shop and its \role{Shop Owner} login are created, and the shop
        app becomes available to that owner at the \texttt{/shop} address.
\end{enumerate}

\subsection{Managing businesses}
From \menu{Platform \textrightarrow{} Businesses} you can:
\begin{itemize}
  \item View every shop on the platform and open a shop's detail page.
  \item See a shop's subscription status, plan, and activity.
  \item Suspend or re-activate a shop.
  \item Adjust a shop's plan or extend its subscription.
\end{itemize}

\subsection{Plans and subscriptions}\label{sec:plans}
\begin{itemize}
  \item \menu{Platform \textrightarrow{} Plans} --- create and edit the
        subscription plans shops can be placed on (name, price, limits).
  \item \menu{Platform \textrightarrow{} Subscriptions} --- see which shop is on
        which plan, renewal dates, and status.
  \item \menu{Platform \textrightarrow{} Payments} --- record and review
        subscription payments.
\end{itemize}

\begin{note}
A shop must have an active subscription to \emph{post} documents (invoices,
purchases). Reading and browsing is always allowed; only posting is gated. If a
shop owner reports a ``subscription required'' message while saving an invoice,
check the shop's subscription here.
\end{note}

\subsection{Users, roles and audit}
\begin{itemize}
  \item \menu{Admin \textrightarrow{} Users} --- platform-level users
        (administrators and sales staff).
  \item \menu{Admin \textrightarrow{} Roles} --- what each platform role may do.
  \item \menu{Platform \textrightarrow{} Sales team} --- manage sales staff who
        onboard and support shops.
  \item \menu{Admin \textrightarrow{} Audit logs} --- a record of important
        actions taken on the platform.
  \item \menu{Admin \textrightarrow{} Settings} --- platform-wide settings.
\end{itemize}

% ===========================================================================
\clearpage
\section{Part B --- Shop Owner \& Staff Guide}

This section is for the people who run a shop day to day: the \role{Shop Owner}
(role: Admin) and \role{Staff} (cashiers, store, accounts).

\subsection{Opening the shop app}
Go to \url{https://apnahisab.sarathitechnologies.com/shop/login} and sign in with
the credentials your administrator gave you. You land on the \menu{Dashboard}.

\subsubsection{Finding your way around}
\begin{itemize}
  \item The \textbf{sidebar} on the left groups every screen (Sales, Purchases,
        Stock, Parties, Money, Reports, and so on). It collapses on small
        screens into a bottom tab bar.
  \item Press \key{Ctrl}+\key{K} (or \key{Cmd}+\key{K} on Mac) anywhere to open
        the \textbf{command menu} and jump straight to any screen.
  \item Most list screens have search, date-range filters, and a
        \menu{Print} button.
\end{itemize}

\subsection{Dashboard}
The home screen summarises today's business --- sales, collections, dues, low
stock and recent activity. Cashiers see a simplified ``My Day'' view focused on
their own sales and receipts.

\subsection{Billing a sale}
There are two ways to bill:
\begin{enumerate}
  \item \menu{Quick Billing} --- the fastest counter screen: pick products, set
        quantities, take payment, done. Best for over-the-counter sales.
  \item \menu{Sales \textrightarrow{} New} --- the full invoice form with every
        option (customer, GST, discounts, line-level details).
\end{enumerate}
When you save (post) an invoice, the system automatically reduces stock, records
the customer's due (if on credit), and updates your accounts.

\subsubsection{After a sale}
\begin{itemize}
  \item A review screen appears; from there you can open the \menu{Share} dialog
        to send the invoice to the customer by a link.
  \item Find all invoices under \menu{Sales \textrightarrow{} Invoices}.
  \item Handle a return under \menu{Sales \textrightarrow{} Returns}.
\end{itemize}

\subsection{Purchases}
\begin{itemize}
  \item \menu{Purchases \textrightarrow{} New} --- record a purchase bill from a
        supplier. Posting it increases stock (at weighted-average cost) and
        records what you owe the supplier.
  \item \menu{Purchases \textrightarrow{} Payables} --- see and pay supplier dues.
  \item \menu{Purchases \textrightarrow{} Returns} --- return goods to a supplier.
\end{itemize}

\subsubsection{AI bill import}
\menu{Purchases \textrightarrow{} Bills} lets you \textbf{upload a photo or PDF of
a supplier bill} and the system reads it for you --- supplier, items, quantities,
taxes --- and prepares a draft. Review the draft (it flags anything uncertain),
fix any line, then confirm to post it. One file can contain several bills.

\subsection{Stock \& inventory}
\begin{itemize}
  \item \menu{Stock \textrightarrow{} Products} --- your product list; open a
        product for its full history and current stock.
  \item \menu{Stock \textrightarrow{} In} / \menu{Out} --- manual stock movements.
  \item \menu{Stock \textrightarrow{} Adjustments} --- correct stock after a count.
  \item \menu{Stock \textrightarrow{} Low stock} --- items running low.
  \item \textbf{Storage locations} --- tag where each item is kept (rack, shelf,
        fridge). Set a product's usual spot, and pick a rack per line while
        billing or purchasing.
  \item \textbf{Combos} --- sell a bundle (e.g. a gift pack) as one product; its
        contents leave stock automatically when sold.
  \item \textbf{Offers} --- set up deals such as buy-2-get-1; free lines are added
        to the sale automatically.
\end{itemize}

\subsection{Parties: customers \& suppliers}
A \textbf{party} is the person or business behind a customer and/or supplier.
\begin{itemize}
  \item \menu{Parties} --- the combined master; a single party can be both a
        customer and a supplier.
  \item \menu{Customers} --- add, import (bulk), and open a customer for their
        ledger and dues.
  \item \menu{Suppliers} --- the same for suppliers.
  \item \menu{Credit Book} --- outstanding receivables and payables, statements,
        and collection reminders.
\end{itemize}

\begin{note}
Money you are owed (receivable) and money you owe (payable) are always shown
separately and are never cancelled out against each other.
\end{note}

\subsection{Money}
\begin{itemize}
  \item \menu{Money \textrightarrow{} Received} --- record a payment from a
        customer; it is allocated against their dues.
  \item \menu{Money \textrightarrow{} Paid} --- record a payment to a supplier.
  \item \menu{Cash \& Bank} --- your cash, bank and UPI accounts, with transfers
        between them.
  \item \menu{Expenses} --- record shop expenses; a single expense can be split
        across payment methods.
\end{itemize}

\begin{note}
Every amount on screen comes straight from the system's records. The app never
does its own arithmetic in the browser, so what you see is what is posted.
\end{note}

\subsection{GST}
If GST is enabled for your shop, GST fields appear on invoices and the
\menu{GST} section gives you GSTR-1 and GSTR-3B summaries and reconciliation.
GST options are hidden if your shop is not registered for GST.

\subsection{Reports}
\menu{Reports} covers Sales, Purchases, Stock, Customers, Suppliers, Expenses,
Profit \& Loss, and GST. Each report supports date ranges and printing.

\subsection{Settings \& staff}
\begin{itemize}
  \item \menu{Settings} --- your shop's profile, preferences, money accounts and
        storage locations.
  \item \menu{Settings \textrightarrow{} Staff} --- add staff and set what each
        one can do (see \S\ref{sec:roles}).
  \item \menu{Profile} --- your own name and password.
  \item \menu{Subscription} --- your current plan and status.
  \item \menu{Recycle Bin} --- recover items you removed.
\end{itemize}

% ===========================================================================
\clearpage
\section{Staff Roles \& Permissions}\label{sec:roles}

The \role{Shop Owner} (Admin) can do everything in the shop and is the only one
who can post or cancel certain documents. \role{Staff} are given a role and a set
of permissions.

\begin{center}
\begin{tabularx}{\textwidth}{@{}p{0.22\textwidth}L@{}}
\toprule
\textbf{Role} & \textbf{Typical access} \\
\midrule
\role{Shop Owner (Admin)} &
Full access. Posts and cancels invoices and purchases, manages staff, sees all
reports. \\[4pt]
\role{Cashier} &
Bills sales and takes payments; sees only their \emph{own} sales, receipts and
returns; no access to reports. \\[4pt]
\role{Store} &
Stock and product work --- goods in/out, adjustments, low stock. \\[4pt]
\role{Accounts} &
Payments, expenses, ledgers and reports. \\
\bottomrule
\end{tabularx}
\end{center}

\begin{note}
The menu only shows what a user is allowed to use. If a colleague cannot see a
screen you can, it is because of their role --- the Shop Owner can adjust this
under \menu{Settings \textrightarrow{} Staff}.
\end{note}

% ===========================================================================
\clearpage
\section{Common Tasks --- Quick Reference}

\begin{center}
\begin{tabularx}{\textwidth}{@{}p{0.40\textwidth}L@{}}
\toprule
\textbf{I want to\ldots} & \textbf{Go to} \\
\midrule
Bill a walk-in customer fast & \menu{Quick Billing} \\
Make a detailed GST invoice  & \menu{Sales \textrightarrow{} New} \\
See all my invoices          & \menu{Sales \textrightarrow{} Invoices} \\
Enter a supplier bill        & \menu{Purchases \textrightarrow{} New} \\
Scan a paper bill            & \menu{Purchases \textrightarrow{} Bills} (upload) \\
Pay a supplier               & \menu{Purchases \textrightarrow{} Payables} \\
Record money from a customer & \menu{Money \textrightarrow{} Received} \\
Add / edit a product         & \menu{Stock \textrightarrow{} Products} \\
Check who owes me            & \menu{Credit Book} \\
Add a staff member           & \menu{Settings \textrightarrow{} Staff} \\
See profit for a period      & \menu{Reports \textrightarrow{} Profit \& Loss} \\
Jump anywhere quickly        & Press \key{Ctrl}+\key{K} \\
\bottomrule
\end{tabularx}
\end{center}

% ===========================================================================
\section{Help \& Troubleshooting}

\begin{description}[leftmargin=0pt,style=nextline]
  \item[I can't sign in.] Check the address matches your role --- shops use
        \texttt{/shop/login}, the administrator uses \texttt{/console/admin/login}.
        Confirm your email and password; ask for a reset if needed.
  \item[A page says a feature isn't available.] That feature may not be enabled
        for your shop, or your role may not include it. Ask the Shop Owner.
  \item[``Subscription required'' when saving.] The shop's subscription is
        inactive; the administrator can renew it from the admin panel.
  \item[The admin panel and shop keep logging each other out.] They share one web
        address, so one browser holds one session. Use a second browser or an
        incognito window for the other panel.
  \item[A print looks wrong.] Use the \menu{Print} button on the screen (not the
        browser menu); it formats the page for printing.
\end{description}

\vfill
\begin{center}
{\color{muted}\small End of guide \textbullet{} Apna Hisab \textbullet{} Sarathi Technologies}
\end{center}

\end{document}
